% TOP_PROBLEM: {"problemId": "problem.skorobogatov-brauer-manin-conjecture-for-k3-surfaces", "canonicalTitle": "Brauer-Manin obstruction to the Hasse principle for K3 surfaces", "releaseRank": 383, "primaryDomain": "number_theory_arithmetic_geometry", "primaryDomainLabel": "Number theory & arithmetic geometry", "displayStatus": "open_with_solved_subcases", "releaseStatus": "open_with_solved_subcases"}
% !TeX program = lualatex
% ATTEMPT_STATUS: unresolved
\documentclass[11pt]{article}
\usepackage[margin=25mm]{geometry}
\usepackage{fontspec}
\setmainfont{Latin Modern Roman}
\usepackage{amsmath,amssymb,mathtools}
\usepackage{unicode-math}
\setmathfont{Latin Modern Math}
\usepackage{xurl,hyperref}
\hypersetup{colorlinks=true,urlcolor=blue}
\setcounter{secnumdepth}{0}
\setlength{\parindent}{0pt}
\setlength{\parskip}{5pt}
\setlength{\emergencystretch}{3em}
\begin{document}
\section*{\texorpdfstring{Brauer-Manin obstruction to the Hasse principle for K3 surfaces}{Brauer-Manin obstruction to the Hasse principle for K3 surfaces}}\label{383}
\subsection{Definitions and mathematical statement}
A K3 surface over a number field $k$ is a smooth projective geometrically integral surface $X$ with $K_X\cong\mathcal O_X$ and $H^1(X,\mathcal O_X)=0$. Let ${\mathbb{A}}_k$ be the adele ring and $\operatorname{Br}(X)=H^2_{\mathrm{et}}(X,\mathbf G_m)$. Define
\[
 X({\mathbb{A}}_k)^{\operatorname{Br}(X)}=
 \left\{(x_v):\sum_v\operatorname{inv}_v(\alpha(x_v))=0
 \text{ for every }\alpha\in\operatorname{Br}(X)\right\},
\]
using the local invariant maps $\operatorname{Br}(k_v)\to{\mathbb{Q}}/{\mathbb{Z}}$; the sum has only finitely many nonzero terms for a fixed class and adelic point. The conjecture is that nonemptiness of this set implies $X(k)\ne\varnothing$. One adelic point must annihilate the entire Brauer group simultaneously. Density of rational points, weak approximation, and obstruction by one single Brauer class are not required.
\par\smallskip\textit{Formulation and concept references:} \href{https://arxiv.org/html/2502.08723v1}{[S1]}.

\subsection{Short English statement}
For every K3 surface over a number field, does surviving all Brauer--Manin tests guarantee at least one rational point?
\subsection{Research attempt: simultaneous Brauer tests and the rational-point gap}
\textbf{Status.} The sufficiency of the Brauer--Manin obstruction for every K3 surface over a number field is \textbf{UNRESOLVED} by this attempt. We establish the necessary direction from reciprocity, reduce the simultaneous obstruction test to a finite compatibility problem using a known finiteness theorem, and verify explicit K3 examples with rational points. None of these steps turns local compatibility into a global rational point in the general case.

\textbf{Source audit.} Monnet's paper [S1] states the K3 Hasse-principle assertion as Conjecture 1.1 and studies explicit bounds for the transcendental Brauer group under principal complex multiplication. Those bounds do not prove the conjecture. We also checked the primary Skorobogatov--Zarhin paper [E1], whose theorem gives finiteness of $\operatorname{Br}(X)/\operatorname{Br}_0(X)$ for K3 surfaces over fields finitely generated over $\mathbb Q$, hence over number fields. Finiteness is not a computed list of generators and is not an effectivity result by itself. The pairing convention and continuity facts below are the standard Brauer--Manin setup, with local and global class-field-theory inputs stated explicitly.

\subsubsection{1. Rational points necessarily survive every test}
Let $P\in X(k)$ and let $\alpha\in\operatorname{Br}(X)$. Evaluation is pullback along $P:\operatorname{Spec}k\to X$, so $P^*\alpha\in\operatorname{Br}(k)$. The image of $P$ in every $X(k_v)$ satisfies
\[
 \alpha(P_v)=\operatorname{res}_{k_v/k}(P^*\alpha).
\]
Global reciprocity for the Brauer group of a number field says that a class $\beta\in\operatorname{Br}(k)$ has
\[
 \sum_v\operatorname{inv}_v(\beta|_{k_v})=0.            \tag{1}
\]
This is an external class-field-theory theorem. Applied to $P^*\alpha$, it gives
\[
 X(k)\subseteq X(\mathbb A_k)^{\operatorname{Br}(X)}.    \tag{2}
\]
The same rational point $P$ works for all $\alpha$, so (2) already has the simultaneous quantifier required in the statement.

Write $\operatorname{Br}_0(X)$ for the image of $\operatorname{Br}(k)$. A constant class evaluates at every local point as the same corresponding local field class, independently of that point. Equation (1) therefore shows that constant classes impose no restriction on adelic points. Adding a constant class to $\alpha$ does not change its global Brauer--Manin functional.

This explains why the quotient by $\operatorname{Br}_0(X)$, rather than the entire generally infinite group $\operatorname{Br}(k)$, is the relevant finite object. It does not justify discarding transcendental classes: they can have nonconstant evaluations and are part of the target.

\subsubsection{2. Finitely many generators suffice, but they must be tested jointly}
Use the theorem of Skorobogatov--Zarhin stated in the source audit. Choose classes $\alpha_1,\ldots,\alpha_r$ whose images generate the finite quotient
\[
 B=\operatorname{Br}(X)/\operatorname{Br}_0(X).
\]
For a fixed adelic point $x=(x_v)$, bilinearity and the preceding constant-class observation give
\[
 x\in X(\mathbb A_k)^{\operatorname{Br}(X)}
 \quad\Longleftrightarrow\quad
 \sum_v\operatorname{inv}_v\alpha_j(x_v)=0
 \quad(1\le j\le r).                                 \tag{3}
\]
Indeed any class differs from an integer linear combination of the $\alpha_j$ by a constant class. Conversely the full group includes every generator. Relations among the generators do not invalidate testing them all; they merely make some equations redundant.

Each chosen $\alpha_j$ is torsion. Let $m$ be a common multiple of their orders as classes in $\operatorname{Br}(X)$, not just of the orders of their images in $B$. Then every local invariant lies in
$m^{-1}\mathbb Z/\mathbb Z$. This choice avoids an unjustified inference that a representative has exactly the order of its class modulo constants.

For these finitely many classes there is a finite set $S$ of places, including the archimedean places, outside which their evaluations vanish at all local points. Here are the standard geometric ingredients. Spread $X$ and the finitely many Brauer classes to a proper smooth model after inverting finitely many primes; classes may be represented by Azumaya algebras on the projective smooth variety, and their finitely many data spread out. Properness extends a $k_v$-point to an integral point of this model. Evaluation then lands in the Brauer group of the local integer ring, which is zero for a complete discretely valued field with finite residue field. Enlarge $S$ to cover the finitely many exceptional primes for all representatives. These are external model and local-Brauer facts; no effective exceptional set is computed here.

Thus a single finite set $S$ works for all points and all generators in (3). This is stronger than observing, separately for one adelic point and one class, that its sum is finite.

\subsubsection{3. An exact finite compatibility criterion}
Assume first that $X(k_v)\ne\varnothing$ for every place, so that local solubility is not already the obstruction. For $v\in S$, form the joint evaluation image
\[
 E_v=\left\{\bigl(\operatorname{inv}_v\alpha_1(P),\ldots,
                  \operatorname{inv}_v\alpha_r(P)\bigr):P\in X(k_v)\right\}
 \subseteq G=(m^{-1}\mathbb Z/\mathbb Z)^r.             \tag{4}
\]
This is a finite subset of a finite abelian group. It need not be a subgroup, and its coordinates need not vary independently. The simultaneous Brauer--Manin condition is exactly
\[
 0\in \sum_{v\in S}E_v.                               \tag{5}
\]
To prove the forward implication, take the evaluation vectors of a Brauer--Manin adelic point and use (3). Conversely, if vectors in (4) sum to zero, select a local point realizing each chosen vector for $v\in S$ and arbitrary local points elsewhere. Since $X$ is proper, its adelic point set is the product of its local point sets; the resulting tuple is adelic. All evaluations outside $S$ vanish, so (3) proves the converse.

Given the exact finite sets $E_v$, (5) is a finite dynamic program. Start with $A_0=\{0\}$ and successively replace
$A_{i+1}=\{a+e:a\in A_i,e\in E_{v_{i+1}}\}$. Induction shows $A_i$ is exactly the set of sums over the first $i$ places. The final membership test is exact and uses at most $|G|$ stored sums at each stage. This is a conditional algorithm: supplying the full generators, their orders, the common exceptional set and all joint local images is part of the input, and may be difficult.

Even perfect execution of this algorithm only determines whether the obstruction vanishes. The conjecture is precisely the further implication from (5) to existence of a $k$-rational point. Counting the finite states does not construct such a point.

\subsubsection{4. A quantifier stress test with all classes, not just generators}
It is easy to confuse
\[
 \forall\alpha\ \exists x\text{ annihilating }\alpha
 \quad\text{with}\quad
 \exists x\ \forall\alpha\text{ annihilating }\alpha.
\]
The distinction persists even for a finite group and even if every linear combination of generators is individually tested. Consider the abstract evaluation set
\[
 E=\{(1,0),(0,1),(1,1)\}\subset\mathbb F_2^2.
\]
For every linear functional $\lambda:\mathbb F_2^2\to\mathbb F_2$, there is $e\in E$ with $\lambda(e)=0$: for a nonzero functional its kernel contains exactly one nonzero vector, while the zero functional vanishes on all three. Nevertheless $0\notin E$, so there is no vector annihilated by both coordinate functionals simultaneously.

Interpreting $0,1$ as $0,1/2$ in $\mathbb Q/\mathbb Z$ gives the same finite evaluation algebra as an exponent-two Brauer test. This is a logical model, not an asserted realization by a K3 surface. It proves that separate success for each class does not formally imply success for the group, so an actual arithmetic proof must preserve the joint images in (4).

A weaker error would replace $E_v$ by the Cartesian product of its coordinate projections. That can insert nonexistent local evaluation combinations. For example the set $\{(0,1),(1,0)\}$ has both coordinate projections equal to $\mathbb F_2$, yet does not contain $(0,0)$. A coordinatewise local search would return a false witness for the simultaneous condition.

These finite examples are exhaustively checked in the diagnostic script. They are tests of quantifiers and algorithms, not counterexamples to the Hasse-principle conjecture.

\subsubsection{5. Compactness gives finite witnesses for failure, not rational points}
For a proper variety over a local field, $X(k_v)$ is compact. For a projective variety this follows by embedding into projective space, whose local points are compact, and taking a closed subset. The adelic product is therefore compact by the product compactness theorem.

For each Brauer class $\alpha$, its global evaluation functional on the adelic space is continuous with finite image: only finitely many local places can contribute for that class, and local evaluation is locally constant into the torsion subgroup of $\mathbb Q/\mathbb Z$ (at real places it is constant on connected components). Hence the annihilator set $A_\alpha$ is closed.

If the full intersection $\bigcap_\alpha A_\alpha$ is empty, compactness implies that a finite subcollection already has empty intersection. Indeed the open complements cover the compact adelic space, so finitely many cover it. Conversely, if every finite collection has a common annihilating adelic point, compactness gives a point annihilating all classes. This is the finite intersection property, and its hypothesis concerns finite collections jointly.

For K3 surfaces, the finite quotient theorem makes that conclusion particularly concrete through (3). But compactness supplies only an adelic tuple. The diagonal subset $X(k)$ is not automatically equal to, or dense in, the closed set obtained by intersecting the $A_\alpha$. A proof that replaced this arithmetic step by topological compactness would already assume what needs to be shown.

\subsubsection{6. An explicit smooth K3 with a rational point and a rational line}
Consider over $\mathbb Q$ the quartic surface
\[
 X:\quad x_0^4+x_1^4=x_2^4+x_3^4\subset\mathbb P^3.
                                                               \tag{6}
\]
The four partial derivatives are $4x_0^3,4x_1^3,-4x_2^3,-4x_3^3$; they cannot vanish simultaneously at a projective point in characteristic zero. Thus $X$ is smooth over an algebraic closure. A reducible projective hypersurface of positive-dimensional components would have intersecting components here and be singular along their intersection, so smoothness makes this hypersurface geometrically integral.

Adjunction gives $K_X\cong\mathcal O_X(4-4)\cong\mathcal O_X$. The sequence
$0\to\mathcal O_{\mathbb P^3}(-4)\to\mathcal O_{\mathbb P^3}\to\mathcal O_X\to0$, together with the usual intermediate cohomology vanishing on $\mathbb P^3$, gives $H^1(X,\mathcal O_X)=0$. Thus (6) is a K3 surface with the definition in the catalogue.

It contains the line
\[
 [s:t]\longmapsto[s:t:s:t],
\]
and in particular the rational point $[1:0:1:0]$. Its diagonal adelic image survives every Brauer test by (2), without needing to enumerate the Brauer group. More generally any adelic tuple obtained by choosing local points on this line is orthogonal to every class: the pullback of a class to $\mathbb P^1_\mathbb Q$ is constant, since $\operatorname{Br}(\mathbb P^1_k)=\operatorname{Br}(k)$, and (1) applies. The Brauer-group equality for projective space is a standard external input here.

This example verifies nonvacuity of the selected implication and all of the smoothness and K3 conditions. It proves neither that every K3 has a rational curve over its ground field nor that a geometric rational curve has a ground-field point. Both assertions would be unjustified extensions of this construction.

\subsubsection{7. Extension fields, zero-cycles and search do not bridge the gap}
A point over a finite extension $K/k$ is a closed point of degree $[K:k]$ when its residue field is exactly $K$. It need not descend to $k$. Thus potential density after extending the field, even if available, addresses a different quantifier from $X(k)\ne\varnothing$ for the original field. Brauer groups also change under field extension, so a calculation over $K$ cannot simply be read as the original obstruction test over $k$.

Likewise a zero-cycle of degree one is an integer linear combination of closed points, not necessarily an effective cycle consisting of one rational point. For a purely logical example, the disjoint union of spectra of a quadratic and a cubic field extension has no $k$-point, but the difference of its degree-three and degree-two closed points has degree one. This is not a K3 counterexample; it shows why an extra geometric theorem would be needed to replace a degree-one cycle by a point. An effective degree-one cycle, by contrast, is a rational point because all its positive residue degrees must sum to one.

Enumerating rational projective coordinates is a semidecision procedure for existence: once a rational point is found, it is a certificate. Failing to find one below any finite height bound is not a nonexistence proof. A Brauer--Manin emptiness certificate can prove nonexistence by (2), but when the set is nonempty the enumeration has no termination bound supplied by the present work.

\subsubsection{8. Verification and the first remaining step}
The explicit quartic's smoothness and rational point are symbolic certificates. The finite examples check joint evaluation images and the dynamic-programming equivalence, including all four linear functionals in the exponent-two stress test. The common finite set of relevant places and finite quotient of the Brauer group are used only with their declared external theorems; no generators for an arbitrary input surface have been fabricated.

The strongest reduction is (5), provided its exact arithmetic inputs are known. The first remaining step is an arithmetic existence theorem turning a compatible system of local points annihilating the full Brauer group into a rational point on an arbitrary K3 surface. Neither finiteness, compactness, a bound for the transcendental Brauer group, nor extension-field points gives that implication. The catalogue assertion remains \textbf{UNRESOLVED} by this attempt.

\subsection{Sources}
\begin{itemize}
\item[S1]S. Monnet, Explicit bounds on the transcendental Brauer group of K3 surfaces with principal complex multiplication, arXiv:2502.08723v1 (2025); pairing convention from Tawfik-Newton arXiv:2311.00650v2.
\url{https://arxiv.org/html/2502.08723v1}.
\textit{Formulation source.}
\item[E1]A. N. Skorobogatov and Y. G. Zarhin, \emph{A finiteness theorem for the Brauer group of abelian varieties and K3 surfaces}, Journal of Algebraic Geometry 17 (2008), 481--502; arXiv:math/0605351.
\url{https://arxiv.org/abs/math/0605351}.
\textit{Primary theorem scope checked 27 September 2026: finiteness of the quotient by constant classes for K3 surfaces over finitely generated characteristic-zero fields.}
\end{itemize}
\end{document}
